% ******************************* Phụ lục A ************************************

\chapter{Kế hoạch quản lý dự án}
\label{appendix:pm}

Phụ lục này mô tả cách nhóm tổ chức 15 tuần thực hiện khóa luận: phạm vi công
việc, các mốc, rủi ro, cách bảo đảm chất lượng, phân công, giao tiếp và quản lý
tài liệu, mã nguồn. Đề tài là một nghiên cứu nên kế hoạch xoay quanh thí nghiệm:
mỗi kết luận ở Chương~\ref{chap:results} phải truy ngược được về một lần chạy có
log, seed và cấu hình cố định.

% ------------------------------------------------------------------------------
\section{Phạm vi công việc và ước lượng}

Công việc được chia thành tám gói (Bảng~\ref{tab:pm-wbs}). Thời gian tính theo
tuần làm việc W01--W15 trong khung 15 tuần của đề cương;
mục~\ref{sec:pm-process} trình bày cách các tuần này khớp với lịch thực.

\begin{table}[htbp]
\caption{Phân rã công việc}
\label{tab:pm-wbs}
\centering
\small
\begin{tabular}{p{0.22\textwidth}p{0.40\textwidth}p{0.13\textwidth}p{0.14\textwidth}}
\toprule
\textbf{Gói công việc} & \textbf{Nội dung chính} & \textbf{Tuần} & \textbf{Phụ trách chính} \\
\midrule
1. Khảo sát và định hướng & Đọc paper theo mảng, literature matrix, khoảng trống nghiên cứu, câu hỏi nghiên cứu, giả thuyết & W01--W03 & Cả nhóm \\
2. Dữ liệu & Lấy ViVQA và ViTextVQA từ nguồn chính thức, cố định split, kiểm tra rò rỉ train/test, OCR cache & W02, W05, W08, W14 & SV2 \\
3. Baseline & Tái lập Qwen2.5-VL-3B-Instruct zero-shot (B1), đối chiếu số công bố trên TextVQA (R0), chốt cấu hình suy luận & W04--W07 & SV1 \\
4. Thích nghi bằng QLoRA & Đường ray huấn luyện QLoRA; các mô hình B2, B3, B3-full, C3, adapter dùng chung SH; ablation cấu hình LoRA & W06--W11 & SV3 \\
5. Đánh giá & EM, VQA-Acc, ANLS, F1; kiểm định McNemar ghép cặp; độ ổn định qua seed; bảng đánh giá đầy đủ & W05--W11 & SV4 \\
6. Phân tích & Error analysis, phân loại câu hỏi, hallucination POPE-vi, ngân sách token, ảnh hưởng của chuẩn hóa và tách từ lên metric & W07, W11--W12, W15 & SV4, SV3 \\
7. Prototype & Web demo, hiệu chỉnh ngưỡng có/không, kiểm tra prototype khớp hệ thống đã đánh giá, video demo & W13--W14 & SV4 \\
8. Báo cáo và bảo vệ & Báo cáo tuần, slide review, gói tái lập, khóa luận & Hằng tuần, W15 & Cả nhóm, trưởng nhóm điều phối \\
\bottomrule
\end{tabular}
\end{table}

Ước lượng ban đầu ở W01 giả định có GPU server của trường. Đến W06, nhóm xác
nhận trường không có GPU server, nên chi phí được ước lượng lại theo đơn vị tính
toán (compute unit, CU) của Google Colab. Từ W08, trước mỗi việc cần GPU, nhóm
ước tính số giờ GPU và số CU (L4 khoảng 1,54 CU/giờ, T4 khoảng 1,2 CU/giờ), rồi
đối chiếu với số đọc thực tế trên Colab. Một lần huấn luyện QLoRA đầy đủ trên
ViTextVQA mất khoảng 12 giờ L4 (B2: 11,99 giờ; B3-full: 12,75 giờ), nên số lần
huấn luyện đầy đủ là ràng buộc lớn nhất của kế hoạch.

% ------------------------------------------------------------------------------
\section{Mục tiêu quản lý}

\begin{itemize}
  \item \textbf{Tái lập được:} mọi lần chạy ghi split, seed, phần cứng, siêu
  tham số, thời gian huấn luyện và suy luận.
  \item \textbf{Không chọn kết quả sau khi thấy số:} luật quyết định của mỗi
  phép thử chính được commit lên Git trước khi chạy.
  \item \textbf{Đúng ràng buộc tài nguyên:} chỉ dùng QLoRA, mọi thí nghiệm chạy
  được trên một GPU Colab (T4 hoặc L4).
  \item \textbf{Đúng mốc:} xong nghiên cứu chính trước Review~2 và hội đồng cấp
  khoa.
\end{itemize}

Đến hết W14, thư mục \path{experiments/} có 141 lần chạy, lần nào cũng có log.
Sáu mô hình chính thức được liệt kê trong \path{configs/official_models.json},
file này sinh tự động từ log; bốn adapter lưu trên Google Drive đều khớp mã băm
SHA-256 đã ghi.

% ------------------------------------------------------------------------------
\section{Rủi ro}

Kế hoạch W01 liệt kê tám rủi ro (mô hình quá lớn, thiếu GPU, dữ liệu lỗi,
baseline khó tái lập, OCR kém, huấn luyện không hội tụ, demo khó triển khai,
huấn luyện chậm). Bảng~\ref{tab:pm-risks} ghi lại các rủi ro đã thực sự xảy ra
và cách nhóm xử lý.

\begin{table}[htbp]
\caption{Rủi ro đã xảy ra và cách xử lý}
\label{tab:pm-risks}
\centering
\small
\begin{tabular}{p{0.24\textwidth}p{0.08\textwidth}p{0.10\textwidth}p{0.46\textwidth}}
\toprule
\textbf{Rủi ro} & \textbf{Tuần} & \textbf{Mức} & \textbf{Cách xử lý} \\
\midrule
Trường không có GPU server & W06 & Cao & Chuyển toàn bộ huấn luyện sang Colab; nâng lên Colab Pro ngày 07/10/2026 để dùng L4 \\
fp16 trên T4 làm hỏng câu trả lời & W06 & Cao & Dò kiểu số trên 73 câu: fp16 hỏng 23 câu; chuyển sang bf16 trên L4 (0/73 câu hỏng) \\
Dữ liệu sai nguồn & W05 & Cao & Dữ liệu trên máy là ViVQA-X chứ không phải ViVQA, còn ViTextVQA lấy từ bản sao không chính thức. Tải lại ViVQA từ repo gốc, xin quyền kho gated của ViTextVQA, dùng đáp án \path{test_gt.json} \\
Một lần huấn luyện đầy đủ tốn khoảng 12 giờ & W08 & Trung bình & Chạy phép thử chính trên tập con 8.000 câu (B3 so với B2-sub); chỉ chạy bản đầy đủ B3-full ở W10 khi kết quả tập con đã rõ \\
Thư viện đổi API & W06 & Thấp & \texttt{transformers} 5.17.0 bỏ \path{warmup_ratio}; đổi sang \path{warmup_steps}, ghi phiên bản thư viện vào log \\
PaddleOCR bản GPU xung đột với torch trong prototype & W13 & Trung bình & Chạy OCR trong một tiến trình riêng (\path{src/prototype/ocr_worker.py}) \\
Adapter dùng chung thiên về trả lời ``có'' & W12--W13 & Trung bình & Chọn ngưỡng trên logit token đầu bằng tập calibration tách rời tập POPE-vi \\
Ảnh có cờ xoay EXIF làm sai OCR cache & W14 & Trung bình & Kiểm tra theo luật ghi trước (164/165 ảnh có cờ xoay bị lệch); sửa \path{ocr.py}, OCR lại 167 ảnh, suy luận lại B3-full (EM 53,35 $\rightarrow$ 53,63) \\
Lịch review thay đổi & W10 & Thấp & Review~2 dời sang 29/10/2026; nhóm cập nhật một bộ slide cộng dồn thay vì làm lại \\
\bottomrule
\end{tabular}
\end{table}

% ------------------------------------------------------------------------------
\section{Quy trình dự án}
\label{sec:pm-process}

Nhóm làm theo chu kỳ tuần. Mỗi tuần làm việc đi qua năm bước:

\begin{enumerate}
  \item \textbf{Lập kế hoạch:} chọn câu hỏi của tuần từ kế hoạch tổng và kết quả
  tuần trước, ghi vào README của tuần.
  \item \textbf{Ghi luật trước:} với mỗi phép thử chính, viết luật quyết định
  (ngưỡng, kiểm định, tập dữ liệu) và commit lên Git trước khi chạy.
  \item \textbf{Chạy:} code được viết và tự kiểm tra (self-test) trên máy cá nhân
  rồi đẩy lên GitHub; notebook Colab clone repo và chạy trên GPU; log, kết quả
  lưu trên Google Drive.
  \item \textbf{Kiểm tra:} tính lại số liệu từ file dự đoán gốc
  (\path{predictions.json}) thay vì chép tay, và clone lại repo để xác minh từng
  commit.
  \item \textbf{Đóng tuần:} cập nhật README, Team Report của tuần và bản nháp các
  chương liên quan.
\end{enumerate}

Tuần làm việc W01--W15 được đánh số theo nội dung trong khung của đề cương, không
theo lịch. Nhóm làm nhanh hơn lịch: ngày 02/11/2026 (tuần thực 9), nhóm đã sang
W15. Tuần thực tính từ thứ Hai 07/09/2026 là tuần 1. Bảng~\ref{tab:pm-weeks} đối
chiếu hai cách đếm.

\begin{table}[htbp]
\caption{Tuần làm việc và lịch thực (năm 2026)}
\label{tab:pm-weeks}
\centering
\small
\begin{tabular}{p{0.12\textwidth}p{0.55\textwidth}p{0.22\textwidth}}
\toprule
\textbf{Tuần} & \textbf{Nội dung chính} & \textbf{Thời gian thực} \\
\midrule
W01--W02 & Định hướng, kế hoạch 15 tuần, literature matrix, phân tích dữ liệu & 07--16/09 \\
W03 & Hồ sơ Review~1 & 17/09 (Review~1: 24/09) \\
W04 & Baseline zero-shot đầu tiên; đổi baseline sang Qwen2.5-VL-3B & 18--21/09 \\
W05 & Dữ liệu chính thức; baseline zero-shot trên cả hai tập & 22--27/09 \\
W06 & Chẩn đoán lỗi fp16; đường ray QLoRA & 28--29/09 \\
W07 & Chốt cấu hình trên L4; B1 chính thức; B2; tái lập R0 & 30/09--12/10 \\
W08 & OCR-enhanced prompting; B3 trên tập con & 12--14/10 \\
W09 & Ổn định qua seed; bảng $2\times2$; ViVQA 2 epoch & 15--17/10 \\
W10 & Ablation cấu hình LoRA; B3-full & 19--21/10 \\
W11 & Bảng đánh giá đầy đủ; adapter dùng chung & 21--22/10 \\
W12 & Hallucination POPE-vi; error analysis & 23--24/10 \\
W13 & Prototype web & 26--31/10 \\
W14 & Hoàn thiện code, video demo, gói tái lập & 01--02/11 \\
W15 & Góp ý Review~2, slide bảo vệ, khóa luận & từ 02/11 \\
\bottomrule
\end{tabular}
\end{table}

Các mốc của khoa giữ nguyên theo lịch thực (Bảng~\ref{tab:pm-milestones}). Điểm
khóa luận gồm 50\% quá trình, trong đó nhóm nộp 7 báo cáo cho giảng viên hướng
dẫn, và 50\% bảo vệ.

\begin{table}[htbp]
\caption{Các mốc của khoa}
\label{tab:pm-milestones}
\centering
\small
\begin{tabular}{p{0.18\textwidth}p{0.22\textwidth}p{0.34\textwidth}p{0.15\textwidth}}
\toprule
\textbf{Mốc} & \textbf{Thời gian} & \textbf{Nội dung} & \textbf{Trạng thái} \\
\midrule
Review~1 & 24/09/2026 (tuần thực 3) & Bài toán, khoảng trống, câu hỏi nghiên cứu, kế hoạch; hai giảng viên không phải GVHD nhận xét & Đã qua, 10 góp ý \\
Review~2 & 29/10/2026 (tuần thực 8) & Baseline đã tái lập, phương pháp, kết quả sơ bộ & Đã qua \\
Hội đồng cấp Khoa/Bộ môn & 30/11--06/12/2026 (tuần thực 13) & Ba thành viên, 60 phút mỗi nhóm & Chưa diễn ra \\
Hội đồng cấp trường & 14--20/12/2026 (tuần thực 15) & Năm thành viên & Chưa diễn ra \\
\bottomrule
\end{tabular}
\end{table}

% ------------------------------------------------------------------------------
\section{Quản lý chất lượng}

Trong đề tài này, chất lượng trước hết là độ tin cậy của con số. Nhóm dùng năm
biện pháp:

\begin{itemize}
  \item \textbf{Luật ghi trước:} ngưỡng và cách kết luận cho mỗi phép thử chính
  được commit trước khi có kết quả, ví dụ luật cho adapter dùng chung (commit
  \texttt{324caeb}, W11) và luật hallucination H4 (commit \texttt{9209170}, W12).
  \item \textbf{Kiểm định ghép cặp:} mọi so sánh hai mô hình trên cùng tập test
  dùng McNemar trên từng câu cho EM và khoảng tin cậy 95\% bằng bootstrap ghép
  cặp cho ANLS; kết quả chính được lặp lại với seed 42, 43 và 44.
  \item \textbf{Tính lại từ dữ liệu gốc:} bảng kết quả cuối
  (\path{src/eval/full_eval_table.py}) đọc trực tiếp các file
  \path{predictions.json}, không chép số từ báo cáo tuần.
  \item \textbf{Self-test:} các module chính có chế độ \texttt{-{}-selftest} chạy
  trên CPU, cả Windows lẫn Linux, để bắt lỗi trước khi tốn GPU.
  \item \textbf{Mã băm và chữ ký run:} adapter, OCR cache, file dữ liệu và slide
  được ghi mã SHA-256; mỗi run mang chữ ký cấu hình nên script không ghi tiếp kết
  quả của cấu hình khác vào cùng thư mục.
\end{itemize}

% ------------------------------------------------------------------------------
\section{Kế hoạch đào tạo}

Đề tài cần kiến thức ở bốn mảng khác nhau, nên ở W01 mỗi thành viên đọc ít nhất
ba paper thuộc mảng của mình và trả lời 10 câu hỏi cố định cho mỗi paper, trong
đó bắt buộc có câu ``dùng được gì cho đề tài'' (Bảng~\ref{tab:pm-papers}).

\begin{table}[htbp]
\caption{Paper nền tảng theo mảng}
\label{tab:pm-papers}
\centering
\small
\begin{tabular}{p{0.32\textwidth}p{0.58\textwidth}}
\toprule
\textbf{Mảng} & \textbf{Paper đã đọc} \\
\midrule
SV1: VLM / Baseline & BLIP-2, Florence-2, Qwen2.5-VL \\
SV2: Dataset / OCR & ViVQA, ViTextVQA, TextVQA (LoRRA) \\
SV3: PEFT / Training & LoRA, QLoRA, VL-Adapter, LLaVA-1.5 \\
SV4: Evaluation / Deployment & POPE, VQAv2, Survey Vision-Language Models (TPAMI 2024) \\
\bottomrule
\end{tabular}
\end{table}

Những kỹ năng còn lại nhóm học trong quá trình làm: huấn luyện QLoRA với PEFT và
bitsandbytes, PaddleOCR và VietOCR, kiểm định McNemar, và cách làm việc giữa máy
cá nhân, GitHub, Colab và Google Drive. Sổ thuật ngữ song ngữ
\path{docs/Glossary.md} được cập nhật suốt dự án.

% ------------------------------------------------------------------------------
\section{Sản phẩm bàn giao}

Bảng~\ref{tab:pm-deliverables} liệt kê sản phẩm theo yêu cầu của đề cương, gồm
cả gói tái lập. Trạng thái tính đến ngày 10/11/2026.

\begin{table}[htbp]
\caption{Sản phẩm bàn giao}
\label{tab:pm-deliverables}
\centering
\small
\begin{tabular}{p{0.38\textwidth}p{0.36\textwidth}p{0.16\textwidth}}
\toprule
\textbf{Sản phẩm} & \textbf{Vị trí} & \textbf{Trạng thái} \\
\midrule
Mã nguồn: 46 module Python (dữ liệu, huấn luyện, đánh giá, prototype) & \path{src/} trên GitHub (PLHGNAOH/ViVQA-VLM) & Hoàn thành \\
30 notebook Colab & \path{notebooks/} & Hoàn thành \\
Log 141 lần chạy & \path{experiments/} & Hoàn thành \\
Adapter QLoRA của các mô hình chính thức & Google Drive; mã băm trong \path{configs/official_models.json} & Hoàn thành \\
Gói tái lập: hướng dẫn, phiên bản thư viện, script suy luận & \path{docs/REPRODUCE.md}, \path{requirements-colab.txt}, \path{src/prototype/infer.py} & Hoàn thành \\
Prototype web và video demo (10 phút 28 giây) & \path{src/prototype/}, Google Drive & Hoàn thành \\
Báo cáo tuần W01--W15 & \path{docs/W01/} đến \path{docs/W15/} & W15 đang làm \\
Slide Review~1, Review~2 (tiếng Anh, tiếng Việt) & \path{docs/W03/}, \path{docs/04_slides/} & Hoàn thành; slide bảo vệ đang làm \\
Khóa luận & Overleaf (bản tiếng Việt trước, bản tiếng Anh sau) & Đang viết \\
\bottomrule
\end{tabular}
\end{table}

% ------------------------------------------------------------------------------
\section{Phân công trách nhiệm}

Nhóm có bốn sinh viên, mỗi người phụ trách một mảng kỹ thuật theo đề cương
(Bảng~\ref{tab:pm-roles}). Phạm Lê Huy Hoàng là trưởng nhóm: ngoài mảng
PEFT/Training, còn điều phối tiến độ và tích hợp kết quả của các mảng. Việc liên
mảng, ví dụ B3 cần OCR cache của SV2 và đường ray huấn luyện của SV3, do người
tạo kết quả chính chịu trách nhiệm và trưởng nhóm duyệt. Đóng góp cụ thể của
từng người theo tuần được ghi trong Team Report.

\begin{table}[htbp]
\caption{Phân công}
\label{tab:pm-roles}
\centering
\small
\begin{tabular}{p{0.20\textwidth}p{0.18\textwidth}p{0.32\textwidth}p{0.20\textwidth}}
\toprule
\textbf{Thành viên} & \textbf{Vai trò} & \textbf{Trách nhiệm chính} & \textbf{Thư mục chính} \\
\midrule
Hồ Đắc Toàn (SE183249) & SV1: VLM / Baseline & Hiểu kiến trúc Qwen2.5-VL, tái lập baseline zero-shot, đối chiếu số công bố & \path{src/models/}, \path{src/eval/run_baseline.py} \\
Trần Quốc Huân (SE193642) & SV2: Dataset / OCR & Dữ liệu chính thức, split, kiểm tra rò rỉ, OCR cache bằng PaddleOCR và VietOCR & \path{data/}, \path{src/data/} \\
Phạm Lê Huy Hoàng (SE182670) & SV3: PEFT / Training; trưởng nhóm & Huấn luyện QLoRA, ablation cấu hình LoRA, adapter dùng chung; điều phối, tích hợp & \path{src/adaptation/}, \path{configs/}, \path{adapters/} \\
Văn Quang Trường (SE182458) & SV4: Evaluation / Deployment & Metric, kiểm định thống kê, hallucination, prototype và demo & \path{src/eval/}, \path{src/prototype/} \\
\bottomrule
\end{tabular}
\end{table}

% ------------------------------------------------------------------------------
\section{Giao tiếp trong dự án}

\begin{table}[htbp]
\caption{Kênh giao tiếp}
\label{tab:pm-comm}
\centering
\small
\begin{tabular}{p{0.26\textwidth}p{0.48\textwidth}p{0.16\textwidth}}
\toprule
\textbf{Kênh} & \textbf{Mục đích} & \textbf{Tần suất} \\
\midrule
Họp với giảng viên hướng dẫn & Xác nhận các quyết định lớn, ví dụ chốt baseline Qwen2.5-VL-3B ngày 07/10/2026 & Theo mốc và khi cần \\
Review với giảng viên phản biện & Nhận góp ý từ bên ngoài; sau mỗi review nhóm điền checklist AI\_Capstone\_Review123, gửi email cho GVHD và giảng viên review, tải lên thư mục chung & Hai lần \\
Báo cáo tuần (README và Team Report trong \path{docs/Wxx/}) & Ghi việc đã làm, kết quả, quyết định và việc dời & Mỗi tuần làm việc \\
GitHub & Mã nguồn, log nhỏ, tài liệu; mỗi commit có mô tả & Liên tục \\
Google Drive & Ảnh, adapter, OCR cache, video demo, ghi âm buổi review & Liên tục \\
Email của bộ môn & Thông báo ngày review, chỉ gửi trong tuần review & Theo mốc \\
\bottomrule
\end{tabular}
\end{table}

% ------------------------------------------------------------------------------
\section{Quản lý cấu hình}

\paragraph{Quản lý tài liệu.} Tài liệu chia làm hai loại. Tài liệu theo tuần
nằm trong \path{docs/W01} đến \path{docs/W15}; mỗi thư mục có README làm mục lục
tuần và một Team Report. Tài liệu xuyên suốt nằm trong các thư mục đánh số:
\path{00_admin} (quy ước, checklist), \path{01_proposal} (đề cương),
\path{02_references} (literature matrix, ghi chú đọc paper), \path{03_report}
(bản nháp các chương), \path{04_slides} (slide review). Mỗi tài liệu chỉ có một
bản sống; thư mục tuần dẫn link tới bản đó thay vì sao chép. Tài liệu cũ được
chuyển vào \path{_archive/} thay vì xóa.

\paragraph{Quản lý mã nguồn.} Mã nguồn được quản lý bằng Git trên repo công khai
\path{github.com/PLHGNAOH/ViVQA-VLM}, làm việc trên nhánh \texttt{main}; nhánh
\texttt{ai-context} do GitHub Actions tự sinh để gom báo cáo. Dữ liệu, ảnh và
trọng số adapter không được đưa lên Git (chặn bằng \path{.gitignore}) mà được
lưu trên Google Drive; vị trí và mã băm được ghi trong log và trong
\path{configs/official_models.json}. Tên lần chạy theo quy ước trong
\path{docs/00_admin/CONVENTIONS.md} và chứa tuần, phương pháp, mô hình, tập dữ
liệu, kiểu số, GPU và seed, ví dụ:
\path{W07_B2_qlora_r16_qwen25vl_vitextvqa_ep1_bf16_L4_seed42}. Mô hình gốc được
ghim theo mã revision trên Hugging Face.

\begin{table}[htbp]
\caption{Công cụ và hạ tầng}
\label{tab:pm-tools}
\centering
\small
\begin{tabular}{p{0.26\textwidth}p{0.64\textwidth}}
\toprule
\textbf{Hạng mục} & \textbf{Công cụ} \\
\midrule
Ngôn ngữ và thư viện & Python, PyTorch, HuggingFace Transformers, PEFT, bitsandbytes, \path{qwen_vl_utils} \\
OCR và xử lý tiếng Việt & PaddleOCR, VietOCR; pyvi để tách từ \\
Huấn luyện và đánh giá & Google Colab Pro: NVIDIA L4 cho huấn luyện và suy luận, Tesla T4 cho OCR (W04--W06 dùng T4 bản miễn phí) \\
Máy cá nhân & Windows 11, GPU 4~GB: phát triển, self-test, tiền xử lý, phân tích lỗi \\
Lưu trữ & Google Drive (dữ liệu, adapter, log lớn); GitHub (mã, log nhỏ, tài liệu) \\
Prototype & Gradio \\
Soạn code và tài liệu & Cursor IDE; trợ lý AI (Claude, Gemini) hỗ trợ viết code và soạn tài liệu, mọi số liệu được đối chiếu lại với log \\
Slide và khóa luận & Canva, PowerPoint, Google Docs, Overleaf \\
\bottomrule
\end{tabular}
\end{table}

Ngân sách Colab được theo dõi bằng CU và ghi trong báo cáo tuần. Số CU còn lại
đọc trên Colab: 146,75 (13/10), 86,52 (19/10), 49,3 (21/10), 34,96 (23/10) và
17,38 (31/10). Riêng W10, với ablation cấu hình LoRA và B3-full, dùng khoảng
37~CU.
